\subsection{Cas 1 : patrimoine hors CTO de un million d'euros}
\label{sec:cas1_sans_retrait}

La figure \ref{fig:sim1} présente la différence d'héritage net pour un investissement de 100 000 euros en assurance-vie ou en CTO, dans le cas où le patrimoine hors enveloppes est de 1 million d'euros.

\begin{figure}[h!]
    \centering
    \includegraphics[width=\textwidth]{cas_1.png}
    \caption{Différence relative d'héritage net pour un investissement de 100 000 euros avant 70 ans, avec un patrimoine hors enveloppes de 1 million d'euros.}
    \label{fig:sim1}
\end{figure}

Dans ce scénario, le poids des autres biens entraîne une forte imposition du CTO, soumis à des tranches élevées du barème successoral (30 à 45 \%). Trois observations peuvent être formulées :
\begin{itemize}
    \item l'assurance-vie domine légèrement (environ 2\% de gain) lorsque la durée de détention est courte (moins de 15 ans), même avec des frais de gestion élevés;
    \item le CTO devient progressivement plus avantageux (jusqu'à un gain relatif de 12\%) lorsque la durée est longue ou lorsque le rendement net de frais est faible;
\end{itemize}
Il est important de noter que les zones où l'assurance-vie conserve un avantage correspondent à des situations défavorables pour l'investisseur (rendements faibles ou décès précoce). Inversement, dans les configurations où le CTO est plus intéressant, l'avantage relatif peut atteindre des niveaux significatifs (jusqu'à 12\% de la succession), alors que l'assurance-vie ne procure qu'un gain marginal (au maximum 2 \%).
